\setcounter{tocdepth}{3}

%%%%%%%%%%%%%%%%%%%%%%%%%%%%%%%%%%%% TITULO TOC %%%%%%%%%%%%%%%%%%%%%%%%%%%%%%%%%%%%%%%%

% \setlength{\cftbeforetoctitleskip}{-2.5em}
\setlength{\cftaftertoctitleskip}{7em}
\renewcommand{\cfttoctitlefont}{\hfill\fontfamily{lmr}\selectfont\resizebox{!}{1.2cm}}
\renewcommand{\cftaftertoctitle}{
    \resizebox{!}{1.2cm}{\,\colorbox{rojoCIMAT}{\!\vphantom{I}\!}}
}

%%%%%%%%%%%%%%%%%%%%%%%%%%%%%%%%%%%% SEPARACIÓN %%%%%%%%%%%%%%%%%%%%%%%%%%%%%%%%%%%%%%%%

% Capitulos
\setlength{\cftchapnumwidth}{2em}
\setlength{\cftbeforechapskip}{0.7em} % arriba
\renewcommand{\cftchapafterpnum}{\vspace{0.5em}} % abajo
% Secciones
\setlength{\cftsecnumwidth}{0em}
\setlength{\cftsecindent}{2em}
\setlength{\cftbeforesecskip}{0.2em} % arriba
% Subsecciones
\setlength{\cftsubsecnumwidth}{2em}
\setlength{\cftsubsecindent}{\cftsecindent} % indentado
% Subsubsecciones
\setlength{\cftsubsubsecindent}{\cftsubsecindent + \cftsubsecnumwidth}

%%%%%%%%%%%%%%%%%%%%%%%%%%%%%%%%%%%%% CAPITULOS %%%%%%%%%%%%%%%%%%%%%%%%%%%%%%%%%%%%%%%%

\cftpagenumbersoff{chap}
% Contorno del renglón; el relleno se ajusta al número por separado.
\newcommand{\chaptertocbox}{%
    \makebox[0pt][l]{%
        \tikz[remember picture,overlay]{
            \draw[color=grisCIMAT, line width=1pt]
                ([yshift=-0.35em]0,0)
                rectangle
                ++(\textwidth,1.22em);
        }%
    }%
}
\renewcommand{\cftchapfont}{\large\chaptertocbox}
\renewcommand{\cftchapaftersnumb}{\scshape}

% Conservamos el ancho mínimo de 1em, la altura y el espacio hasta el título.
\newsavebox{\chaptertocnumber}
\newcommand{\chaptertocnumberline}[1]{%
    \begingroup
    \setlength{\fboxsep}{0pt}%
    \sbox{\chaptertocnumber}{%
        \colorbox{grisCIMAT}{%
            \rule[-0.35em]{0pt}{1.22em}%
            \makebox[\maxof{1em}{\widthof{\bfseries#1}+0.3em}][c]{%
                \bfseries\color{white}#1%
            }%
        }%
    }%
    \makebox[\dimexpr\wd\chaptertocnumber+\cftchapnumwidth-1em\relax][l]{%
        % El fondo no altera el interlineado original de la entrada.
        \smash{\usebox{\chaptertocnumber}}\vphantom{\bfseries#1}%
    }%
    \endgroup
    \cftchapaftersnumb
}

% Aplicamos el número adaptable únicamente a los capítulos, como en los apuntes.
\makeatletter
\let\oldl@chapter\l@chapter
\renewcommand{\l@chapter}[2]{%
    \begingroup
    \let\numberline\chaptertocnumberline
    \oldl@chapter{#1}{#2}%
    \endgroup
}
\makeatother

%%%%%%%%%%%%%%%%%%%%%%%%%%%%%%%%%%%%% SECCIONES %%%%%%%%%%%%%%%%%%%%%%%%%%%%%%%%%%%%%%%%

\renewcommand{\cftsecfont}{\fontfamily{bch}\selectfont\color{grisCIMAT}}
\cftpagenumbersoff{sec}
\renewcommand{\cftsecpresnum}{\color{white}}
\renewcommand{\cftsecdotsep}{\cftnodots}

\renewcommand{\cftsubsecfont}{\small\fontfamily{bch}\selectfont}
\renewcommand{\cftsubsecaftersnum}{.}

\renewcommand{\cftsubsubsecfont}{\small\fontfamily{cmss}\selectfont}
